\section{Einleitung}\label{sec:einleitung}

% Diese Abkürzungen treten im Fließtext ausschließlich als Bestandteil deutscher
% Komposita bzw. in flektierter Form auf und lassen sich daher nicht über \ac{}
% einführen. \acused{} meldet sie dennoch dem Abkürzungsverzeichnis.
\acused{KI}\acused{DORA}\acused{HTTP}\glsadd{devops}%

Große Sprachmodelle sind in der Softwareentwicklung angekommen: 81\,\% der Befragten des DORA-Reports 2024 berichten eine Prioritätenverschiebung zugunsten von KI.\footcite[\vglf][\pagef 18]{DORA.2024} Eine Mehrfallstudie in drei deutschen Organisationen zeigt den regulären Einsatz in agilen Teams bei häufig fehlenden Regelungen.\footcite[\vglf][\pagef 297 f.]{Neumann.2026} Die Modelle treffen dort auf ein Vorgehensmodell aus kurzen Lieferzyklen, technischer Exzellenz und regelmäßiger Reflexion.\footcite[\vglf][o.\,S.]{Beck.2001}

Die Befundlage ist widersprüchlich. Ein Laborexperiment weist eine um 55,8\,\% kürzere Bearbeitungszeit aus, ein randomisierter Versuch mit erfahrenen Entwicklern dagegen eine um 19\,\% längere.\footcites(\vglf)()[][\pagef 2]{Peng.2023}[][\pagef 2]{Becker.2025} Der DORA-Report verbindet zudem eine um 25\,\% höhere KI-Adoption mit besserer Dokumentationsqualität, aber geringerer Auslieferungsstabilität.\footcite[\vglf][\pagef 37, \pagef 39]{DORA.2024} Ob \glspl{codegenerator} agile Prozesse stärken oder aushöhlen, ist damit offen.

Die Arbeit fragt, welche Potenziale und Risiken der Einsatz KI-gestützter Code-Generatoren in agilen Softwareentwicklungsprozessen mit sich bringt und welche Gestaltungsempfehlungen daraus folgen. Ziel ist deren Systematisierung entlang agiler Praktiken.

Methodisch stützt sich die Arbeit auf eine literaturbasierte Studienauswertung; eigene Erhebungen, Werkzeugvergleiche, juristische Detailfragen und KI-Ethik bleiben ausgeklammert. Kapitel~\ref{sec:grundlagen} klärt die Begriffe, Kapitel~\ref{sec:potenziale} und~\ref{sec:risiken} stellen Potenziale und Risiken gegenüber, Kapitel~\ref{sec:wuerdigung} wägt ab, Kapitel~\ref{sec:fazit} beantwortet die Forschungsfrage.
